\documentclass[10pt,journal,compsoc]{IEEEtran}

\usepackage{amsmath,amsfonts,amssymb}
\usepackage{graphicx}
\usepackage{booktabs}
\usepackage{algorithm}
\usepackage{algorithmic}
\usepackage{hyperref}
\usepackage{xcolor}

\begin{document}

\title{TabDrift: Preventing Categorical Collapse in One-Step Tabular Generation via Drifting Vector Fields}

\author{Your Name,~\IEEEmembership{Student Member,~IEEE,}
        and~Your Advisor,~\IEEEmembership{Senior Member,~IEEE}% <-this % stops a space
\IEEEcompsocitemizethanks{\IEEEcompsocthanksitem Y. Name and Y. Advisor are with the Department of Computer Science, University of XYZ, City, State, Zip.\protect\\
E-mail: \{your.email, advisor.email\}@xyz.edu}}

\IEEEtitleabstractindextext{%
\begin{abstract}
Generative modeling for mixed-type tabular data has seen a paradigm shift from Generative Adversarial Networks (GANs) to Diffusion Models, achieving state-of-the-art statistical fidelity. However, this fidelity comes at a prohibitive computational cost: tabular diffusion models require solving probability flow ordinary differential equations (ODEs) over 50 to 100 sequential steps, rendering them impractical for real-time applications or massive-scale synthetic database generation. While 1-step generation methods such as Consistency Distillation have accelerated image synthesis, their application to the discrete-continuous topology of tabular latent spaces often results in severe categorical mode collapse due to high-curvature ODE trajectories. In this paper, we introduce \textbf{TabDrift}, the first adaptation of Drifting Models to the tabular domain. Rather than distilling an ODE, TabDrift learns a native 1-step mapping from noise to data by evolving the pushforward distribution during training. Guided by an Attraction-Repulsion Mean-Shift kernel, the repulsive forces naturally preserve discrete categorical modes within the continuous VAE latent space. Empirical evaluations on standard benchmarks (Adult, Default, Shoppers, Magic) demonstrate that TabDrift achieves statistical fidelity and Machine Learning Efficacy (MLE) competitive with 50-step diffusion models while delivering a $>50\times$ speedup in inference latency, establishing a new state-of-the-art for fast, high-fidelity tabular synthesis.
\end{abstract}

\begin{IEEEkeywords}
Tabular Data Synthesis, Generative AI, Drifting Models, Diffusion Models, One-Step Generation, Mean-Shift
\end{IEEEkeywords}}

\maketitle
\IEEEdisplaynontitleabstractindextext
\IEEEpeerreviewmaketitle

\section{Introduction}
\IEEEPARstart{T}{he} generation of synthetic tabular data is a critical enabler for privacy-preserving data sharing in finance and healthcare, as well as a powerful tool for augmenting imbalanced datasets in machine learning workflows. Unlike continuous domains such as images or audio, tabular data inherently possesses a heterogeneous topology: it features a mix of continuous numerical variables and discrete, often highly imbalanced, categorical variables.

Historically, Generative Adversarial Networks (GANs) such as CTGAN and TVAE were the standard for tabular synthesis due to their fast, single-step inference capabilities. However, GANs are notoriously difficult to train, suffer from mode collapse, and often fail to accurately capture complex inter-column correlations. Recently, score-based generative models and diffusion models (e.g., TabDDPM, CoDi, TabSyn) have emerged as the new state-of-the-art. By mapping mixed-type data to a continuous latent space and slowly reversing a diffusion process, these models achieve unprecedented statistical fidelity. 

Despite their high fidelity, diffusion models are computationally paralyzed during inference. Generating a single synthetic row requires 50 to 100 sequential evaluations of a neural network through an Ordinary Differential Equation (ODE) solver. For enterprise applications requiring millions of synthetic records, this latency is a fundamental barrier to deployment. 

Recent efforts to accelerate diffusion models, such as Consistency Models and Latent Consistency Distillation, have seen immense success in the image domain. However, we hypothesize and demonstrate that directly applying Consistency Distillation to tabular latent spaces is suboptimal. The continuous latent representations of discrete categorical variables create sharp, jagged boundaries in the latent space. Consequently, the Probability Flow ODE trajectories exhibit high curvature. When a 1-step Consistency Model attempts to "cut the corner" of these trajectories, it outputs averaged embeddings, resulting in severe categorical mode collapse.

To bridge the gap between the fidelity of diffusion models and the speed of GANs without suffering from categorical collapse, we propose \textbf{TabDrift}. Inspired by the recently introduced paradigm of Generative Modeling via Drifting, TabDrift trains a single-step generator natively. Rather than matching a pre-trained ODE trajectory, TabDrift evolves its pushforward distribution during training using a mean-shift vector field. Crucially, this drifting field incorporates an explicit \textit{Repulsion} term between generated samples. In the context of tabular latent spaces, this repulsion naturally pushes generated samples away from the "average" center, forcing them into distinct categorical clusters and perfectly preserving discrete modes.

\textbf{Our main contributions are as follows:}
\begin{itemize}
    \item We identify the architectural limitations of applying Consistency Distillation to the high-curvature latent spaces of mixed-type tabular data.
    \item We propose TabDrift, adapting Generative Modeling via Drifting to tabular data, achieving native 1-step generation without the need for pre-training a teacher diffusion model.
    \item We demonstrate that the inherent Attraction-Repulsion kernel of TabDrift acts as a topological regularizer, preventing categorical mode collapse.
    \item We empirically validate that TabDrift reduces inference latency by $>50\times$ compared to state-of-the-art tabular diffusion models while maintaining comparable Machine Learning Efficacy (MLE).
\end{itemize}

\section{Related Work}
\subsection{Tabular Generative Models}
Prior to the advent of diffusion models, GAN-based methods dominated tabular synthesis. CTGAN introduced mode-specific normalization to handle numerical distributions, while TVAE utilized variational lower bounds. While fast, these models often fail to capture joint distributions effectively. Autoregressive models like REaLTabFormer treat tabular rows as text sequences but suffer from sequence-length scaling issues and slow step-by-step token generation. 

\subsection{Tabular Diffusion Models}
TabDDPM applied standard denoising diffusion probabilistic models directly to tabular inputs using a multinomial diffusion process for categorical variables. TabSyn improved upon this by training a continuous VAE to encode mixed-type data into a shared continuous latent space, subsequently applying Elucidated Diffusion Models (EDM) on the latents. While TabSyn achieves SOTA fidelity, it requires $N \approx 50$ ODE solver steps for generation.

\subsection{Fast Generative Models and Drifting}
Consistency Models enable fast generation by learning to map any point on a probability flow ODE to its origin. However, distillation suffers when trajectory curvature is high. Very recently, Deng et al. (2025) proposed Generative Modeling via Drifting, formulating generation as a pushforward evolution driven by a mean-shift vector field. To the best of our knowledge, TabDrift is the first application of drifting models to tabular data.

\section{Methodology}
\subsection{Continuous Latent Space Representation}
Following TabSyn, we isolate the discrete nature of tabular data by employing a pre-trained Variational Autoencoder (VAE). Let $X = [X_{num}, X_{cat}]$ be a tabular dataset containing numerical and categorical columns. The VAE encoder $E_\phi$ maps a row $x \in X$ to a continuous latent representation $z = E_\phi(x) \in \mathbb{R}^D$. Generating synthetic data is therefore reduced to generating synthetic latent vectors $\hat{z}$, which are then decoded via $D_\theta(\hat{z}) \to \hat{x}$.

\subsection{The TabDrift Generator}
Unlike diffusion models which require a time-conditioned network $f_\theta(z_t, t)$, TabDrift utilizes an unconditioned Multi-Layer Perceptron (MLP) generator $G_\psi$. The generator takes standard Gaussian noise $\epsilon \sim \mathcal{N}(0, I)$ and maps it directly to the latent space in a single forward pass: $\hat{z} = G_\psi(\epsilon)$.

\subsection{Topology-Aware Drifting Field}
The core of TabDrift lies in its training objective. At each training step, we sample a batch of real latent vectors $Z_{real}$ and generate a batch of fake vectors $Z_{fake} = G_\psi(\epsilon)$. We compute a drifting vector field $V(\hat{z})$ that dictates how each generated sample should move to better match the real distribution.

The field is defined by two forces:
\begin{equation}
    V(\hat{z}) = V^+(\hat{z}, Z_{real}) - V^-(\hat{z}, Z_{fake})
\end{equation}

Where $V^+$ is the attraction toward real samples, and $V^-$ is the repulsion from other generated samples. The weights are determined by a doubly-stochastic normalized kernel $k(x, y) = \exp(-||x-y|| / (\tau \sqrt{D}))$.

\textbf{Preventing Categorical Collapse:} In the VAE latent space, distinct categorical values (e.g., "Married" vs "Divorced") form distant, dense clusters. If a model generates a sample precisely between these clusters, a standard MSE loss would pull it toward both, leaving it stuck in the "average" (invalid) space. In TabDrift, if multiple fake samples occupy this "average" space, the repulsion term $V^-$ violently pushes them apart, naturally driving them toward the distinct discrete modes of the real data.

\subsection{Optimization}
We freeze the drifting field target and train the generator $G_\psi$ using a simple regression objective:
\begin{equation}
    \mathcal{L} = \mathbb{E}_{\epsilon} \left[ || G_\psi(\epsilon) - \text{stopgrad}(G_\psi(\epsilon) + V(G_\psi(\epsilon))) ||^2 \right]
\end{equation}
This forces the generator to evolve its pushforward distribution until the field reaches equilibrium ($V=0$).

\section{Experiments}
[This section will be populated with the quantitative results from the background training process, comparing TabDrift (1-step), TabSyn (50-steps), and CTGAN (1-step) on Adult, Default, Shoppers, and Magic datasets.]

\section{Conclusion}
TabDrift successfully bridges the gap between the inference speed of GANs and the statistical fidelity of Diffusion models. By leveraging the repulsive properties of drifting vector fields, we solve the categorical mode collapse problem inherent to 1-step tabular generators, enabling massive-scale synthetic database generation in seconds rather than hours.

\end{document}
